\section{Accessions de la cohorte de surveillance}\label{sec:accessions}
Le tableau~\ref{tab:accessions} identifie les séquences publiques utilisées pour construire la matrice de distances. Les 64 enregistrements sont déclarés complets par leurs déposants, selon le filtre \texttt{--complete-only} de NCBI Datasets~18.41.1~\cite{ncbidatasets}.\footnote{Voir la \href{https://ncbi.nlm.nih.gov/datasets/docs/v2/reference-docs/command-line/datasets/download/virus/genome/datasets_download_virus_genome_accession/}{documentation des filtres de NCBI Datasets} et le manifeste \reposource{results/ncbi/sars_metadata.tsv}.} La référence Wuhan-Hu-1 appartient au lignage B. Pour les sept groupes de neuf, le filtre \texttt{--lineage} est suivi d'une vérification de l'étiquette Pango de chaque accession.

\begin{table}[htbp]
\centering
\caption{Les accessions versionnées sont regroupées par lignage Pango. La colonne \(\nu\) donne le nombre de caractères \texttt{N} par accession.}\label{tab:accessions}
\small
\begin{tabular}{@{}llr@{\hspace{1.6em}}lr@{\hspace{1.6em}}lr@{}}\toprule
Groupe & Accession & \(\nu\) & Accession & \(\nu\) & Accession & \(\nu\) \\\midrule
Référence & \texttt{NC\_045512.2} & 0 & & & & \\
A & \texttt{OX620412.1} & 1\,983 & \texttt{OX621768.1} & 246 & \texttt{OX640170.1} & 25 \\
 & \texttt{OX641140.1} & 634 & \texttt{OX654322.1} & 6\,810 & \texttt{OX654867.1} & 1\,045 \\
 & \texttt{OX655642.1} & 0 & \texttt{OX659521.1} & 8\,705 & \texttt{OX659701.1} & 273 \\\addlinespace[0.3ex]
B.1 & \texttt{PX058549.1} & 3 & \texttt{PX058550.1} & 16 & \texttt{PX058551.1} & 17 \\
 & \texttt{PX058553.1} & 0 & \texttt{PX058554.1} & 0 & \texttt{PX058556.1} & 0 \\
 & \texttt{PX058558.1} & 0 & \texttt{PX058560.1} & 0 & \texttt{PX058562.1} & 30 \\\addlinespace[0.3ex]
B.1.1.7 & \texttt{OZ293821.1} & 121 & \texttt{OZ293826.1} & 121 & \texttt{OZ293831.1} & 123 \\
 & \texttt{OZ293860.1} & 121 & \texttt{OZ293864.1} & 121 & \texttt{OZ293867.1} & 2\,174 \\
 & \texttt{OZ293868.1} & 122 & \texttt{OZ293878.1} & 378 & \texttt{OZ293879.1} & 121 \\\addlinespace[0.3ex]
B.1.351 & \texttt{PP520277.1} & 52 & \texttt{PP520338.1} & 36 & \texttt{PP520436.1} & 19 \\
 & \texttt{PP520490.1} & 20 & \texttt{PP520499.1} & 20 & \texttt{PP520501.1} & 20 \\
 & \texttt{PP520502.1} & 27 & \texttt{PP520576.1} & 30 & \texttt{PP520579.1} & 20 \\\addlinespace[0.3ex]
P.1 & \texttt{PQ094852.1} & 0 & \texttt{PQ094853.1} & 0 & \texttt{PQ094855.1} & 0 \\
 & \texttt{PQ094856.1} & 0 & \texttt{PQ094857.1} & 0 & \texttt{PQ094858.1} & 0 \\
 & \texttt{PQ094859.1} & 0 & \texttt{PQ094861.1} & 0 & \texttt{PQ094862.1} & 0 \\\addlinespace[0.3ex]
B.1.617.2 & \texttt{PQ992576.1} & 108 & \texttt{PQ992584.1} & 35 & \texttt{PQ992589.1} & 87 \\
 & \texttt{PQ992594.1} & 48 & \texttt{PQ992629.1} & 52 & \texttt{PQ992650.1} & 14 \\
 & \texttt{PQ992659.1} & 7 & \texttt{PQ992672.1} & 28 & \texttt{PQ992676.1} & 15 \\\addlinespace[0.3ex]
B.1.1.529 & \texttt{OY697980.1} & 174 & \texttt{OY697992.1} & 2\,298 & \texttt{OY698214.1} & 6\,866 \\
 & \texttt{OY698293.1} & 7\,473 & \texttt{OY698556.1} & 6\,577 & \texttt{OY698631.1} & 2\,107 \\
 & \texttt{OY698922.1} & 8\,168 & \texttt{OY699229.1} & 7\,598 & \texttt{OY699620.1} & 477 \\\addlinespace[0.3ex]
\bottomrule
\end{tabular}
\end{table}

\FloatBarrier
Les séquences comptent de 29\,580 à 29\,903 caractères. Sur 1\,905\,372 caractères, 65\,565 sont des \texttt{N} et 23 portent une autre ambiguïté IUPAC. Le coût de substitution est \(c(a,b)=\mathbf{1}_{a\ne b}\), y compris pour \texttt{N}. Ainsi, \texttt{N/N} coûte zéro et \texttt{N/A} coûte un. Si \(\nu(x)\) désigne le nombre de \texttt{N} de la chaîne \(x\), une édition peut changer ce nombre d'au plus une unité. Il en résulte
\[
d(x,y)\;\geq\;|\nu(x)-\nu(y)|.
\]
Les effectifs par accession et les empreintes des séquences figurent dans \reposource{results/ncbi/sars_content.tsv}.

\synthese{Sur 915 des 2\,016 paires, le seul écart entre nombres de \texttt{N} impose au moins la moitié de la distance observée. Pour la paire la plus distante, cette borne vaut 537 éditions sur 12\,229.}
